% Pursuing Stacks (cotes 134-2 to 134-8): archive cote 134-7, site page 53 (PDF page 54), typescript page 495.
% typescript. Draft compared: ps6.tex:2617-2677.
% AI-transcribed by Claude Opus 5.5 (claude-opus-5-5, adaptive thinking, effort high), Message Batches API, batch ps4 (msgbatch_012Aoc1NuskATPsyBeKSYK7J), 2026-11-12; not yet checked by eye.
\scanpage{134-7/53}
\mspage{495}

may be equally viewed as being \uline{pro}-\uline{modules} endowed with suitable extra
structure.) Correspondingly, we will expect the cohomology invariants
constructed in terms of (the apparently more forbidding) $\hand{\Gamma}_k{}^{\wedge}$ to express
quasi-coherent cohomology of the corresponding \uline{schemes}, or semisimplicial
systems of such; whereas, working with the (apparently more anodine !)
$\hand{\Gamma}_k$ will lead to cohomology invariants of \uline{formal schemes} and semisimplici-%
al systems of such. Both types of invariants are of interest it would
seem, the one however which looks the more relevant in connection with
studying ordinary homotopy types in terms of schematic one's, is surely
the first. On the other hand, there is'nt any reason whatever to
believe that under fairly general conditions, these two types of invari-%
ants are going to be isomorphic, by the evident map from ``schematic'' to
``formal'' quasicoherent cohomology. To say it differently, I do no longer
expect that under\unsure{reasonably| reasonably} wide assumptions, the map (24) of p.488
\[
(6)\qquad L_{pt}(X_{*})\unsure{\hand{\longrightarrow}|\longrightarrow} L(X_{*})
\]
is a quasi-isomorphism, nor behaves like one with respect to taking
$\mathrm{Ext}^{i}$ 's with values in a quasicoherent module, as I was \inserted{hastily} surmising for
about one week, while loosing track of the geometric meaning of the
algebraic objects I was playing around with. To give just one example,
take $X_{*}$ to be the standard semisimplicial \xout{object} unipotent bundle associa-%
ted to the group-object $\underline{G}_a$ in $\hand{\mathcal{U}}(k)$, namely just the usual affine line
with addition law. The Ext$^{i}$'s of the two members of (6) with values in
the k-module k may be interpreted as either schematic or formal Hochschild
cohomology of the additive group, with coefficients in \xout{\illegible{1}} $\underline{O}_k = \underline{G}_a$ .
The map of the former into the latter is not \inserted{always} an isomorphism, already
in dimension 1 , where the two groups to be compared are just the groups
of endomorphisms of $\underline{G}_a$ , and of the corresponding formal group. If k is
of char. p $\unsure{\hand{>}|>}$ 0, a prime, then the latter group can be described as the
group of all formal power series of the type
\[
F(t) = \hand{\sum}_{i\geq 0} c_i t^{(p^{i})} ,
\]
whereas for the first group we must restrict to those F which are polyno-%
mials i.e. \xout{all} only a finite number of the coefficients $c_i$ are non-zero.

One may of course object to this example, because the $X_{*}$ we are
working with is not simply connected, and \deleted{that} \inserted{because} the example does not apply
\deleted{either} over a ring such as $\underline{Z}$ , which is the one we are interested in most
of all. I am convinced \inserted{now} however that even when assuming $X_1 = X_0 = e$
and k=$\underline{Z}$ , (6) is very far from being a quasiisomorphism, even for such

\margin{53}
